\section*{Bab \ref{ch:g11:seq} --- Barisan: Perkenalan Pertama}

\begin{solution}{exo:g11:seq:1}
$u_n = \frac{n}{n+1}$: $u_1 = \frac12$, $u_2 = \frac23$,
$u_3 = \frac34$.

$u_0 = 5$, $u_{n+1} = 3u_n - 2$: $u_1 = 13$, $u_2 = 37$, $u_3 = 109$.

$u_n = (-1)^n n$: $u_1 = -1$, $u_2 = 2$, $u_3 = -3$.
\end{solution}

\begin{solution}{exo:g11:seq:2}
$u_{10} = 7 + 10 \times (-3) = -23$ dan $u_{25} = 7 - 75 = -68$.

Untuk $(v_n)$: $v_8 = v_3 + 5d$ memberi $26 = 11 + 5d$, jadi $d = 3$; lalu
$v_0 = v_3 - 3d = 11 - 9 = 2$.
\end{solution}

\begin{solution}{exo:g11:seq:3}
$u_8 = 5 \times 2^8 = 1280$.

Untuk $(v_n)$: $v_4 = v_2\, q^2$ memberi $48 = 12 q^2$, jadi $q^2 = 4$ dan
$q = 2$ (sukunya positif). Lalu
$v_0 = \frac{v_2}{q^2} = \frac{12}{4} = 3$.
\end{solution}

\begin{solution}{exo:g11:seq:4}
$1 + \dots + 500 = \frac{500 \times 501}{2} = 125\,250$.

$4 + 7 + \dots + 61$ bersifat \omterm{def:g11:seq:arithmetic}{aritmetika} dengan $d = 3$ dan
$\frac{61 - 4}{3} + 1 = 20$ suku: jumlahnya
$20 \times \frac{4 + 61}{2} = 650$.

$1 + \frac12 + \dots + \frac{1}{2^{10}}$ bersifat \omterm{def:g11:seq:geometric}{geometri} dengan
$q = \frac12$ dan $11$ suku:
$\frac{1 - (1/2)^{11}}{1 - 1/2} = 2\left(1 - \frac{1}{2048}\right)
= \frac{2047}{1024}$.
\end{solution}

\begin{solution}{exo:g11:seq:5}
$u_{n+1} - u_n = 4(n+1) - 1 - 4n + 1 = 4$: \omterm{def:g11:seq:arithmetic}{aritmetika} dengan $d = 4$.

$\frac{v_{n+1}}{v_n} = \frac{2^{n+1}}{3^{n+2}} \cdot \frac{3^{n+1}}{2^n}
= \frac23$: \omterm{def:g11:seq:geometric}{geometri} dengan $q = \frac23$.

$w_0 = 0$, $w_1 = 2$, $w_2 = 6$: selisihnya $2$ dan $4$ berbeda, jadi
bukan \omterm{def:g11:seq:arithmetic}{aritmetika}; $\frac{w_1}{w_0}$ bahkan tidak terdefinisi, dan \omterm{def:g11:seq:geometric}{rasionya}
$\frac{w_2}{w_1} = 3 \neq \frac{w_3}{w_2} = 2$: jadi bukan keduanya.
\end{solution}

\begin{solution}{exo:g11:seq:6}
Banyaknya kursi per baris bersifat \omterm{def:g11:seq:arithmetic}{aritmetika}: suku pertama $16$, \omterm{def:g11:seq:arithmetic}{bedanya} $2$.
Baris terakhir (baris ke-20) mempunyai $16 + 19 \times 2 = 54$ kursi. Seluruhnya
$20 \times \frac{16 + 54}{2} = 700$ kursi.
\end{solution}

\begin{solution}{exo:g11:seq:7}
Setelah $n$ jam, populasinya $500 \times 2^n$. Pada pukul 20.00, $n = 8$:
$500 \times 256 = 128\,000$ bakteri. Kita perlu $500 \times 2^n >
10^6$, yaitu $2^n > 2000$: karena $2^{10} = 1024$ dan $2^{11} = 2048$,
populasinya pertama kali melampaui satu juta setelah $11$ jam penuh, yaitu
pada pukul 23.00.
\end{solution}

\begin{solution}{exo:g11:seq:8}
$c_2 = 1.002 \times 100 + 100 = 200.20$ dan
$c_3 = 1.002 \times 200.20 + 100 \approx 300.60$. Selisihnya,
$c_2 - c_1 = 100.20$ dan $c_3 - c_2 \approx 100.40$, tidak sama, sehingga
$(c_n)$ bukan \omterm{def:g11:seq:arithmetic}{aritmetika}; \omterm{def:g11:seq:geometric}{rasionya}, $\frac{c_2}{c_1} = 2.002$ dan
$\frac{c_3}{c_2} \approx 1.50$, juga tidak sama, sehingga \omterm{def:g11:seq:sequence}{barisan} itu bukan
\omterm{def:g11:seq:geometric}{geometri}. (Rekursi campuran ``kalikan lalu tambahkan'' seperti ini
diselesaikan dengan kiat \omterm{def:g11:seq:sequence}{barisan} bantu pada \cref{exo:g11:seq:11}.)
\end{solution}

\begin{solution}{exo:g11:seq:9}
$u_{n+1} - u_n = (n+1)^2 - 8(n+1) - n^2 + 8n = 2n - 7$: negatif untuk
$n \leq 3$, positif untuk $n \geq 4$. Jadi $(u_n)$ turun sampai
$u_4 = 16 - 32 = -16$, lalu naik: \omterm{def:g11:seq:sequence}{barisan} itu tidak monoton.

$(v_n)$ mempunyai suku positif dan
\[
\frac{v_{n+1}}{v_n} = \frac{3^{n+1}}{(n+1)!} \cdot \frac{n!}{3^n}
= \frac{3}{n+1},
\]
yang bernilai $> 1$ untuk $n \leq 1$, $= 1$ untuk $n = 2$, dan $< 1$ untuk
$n \geq 3$: jadi \omterm{def:g11:seq:sequence}{barisannya} naik sampai $v_2 = v_3 = \frac92$, lalu
turun.
\end{solution}

\begin{solution}{exo:g11:seq:10}
Sebanyak $n$ suku pertamanya adalah $u_0, \dots, u_{n-1}$, dengan $u_0 = 3$
dan $u_{n-1} = 3 + 4(n-1) = 4n - 1$. Jumlahnya adalah
\[
n \times \frac{3 + (4n-1)}{2} = n(2n + 1) = 903,
\]
sehingga $2n^2 + n - 903 = 0$. Di sini $\Delta = 1 + 4 \times 2 \times 903 =
7225 = 85^2$, dan $n = \frac{-1 + 85}{4} = 21$ (\omterm{def:g11:quad:discriminant}{akar} negatifnya
ditolak). Periksa: $21 \times 43 = 903$.
\end{solution}

\begin{solution}{exo:g11:seq:11}
\emph{1.} $u_1 = 5$, $u_2 = 9$, $u_3 = 17$: setiap sukunya satu lebih besar
daripada $4, 8, 16$, yang mengisyaratkan $u_n = 2^{n+1} + 1$.

\emph{2.} Dengan $v_n = u_n - 1$:
\[
v_{n+1} = u_{n+1} - 1 = 2u_n - 1 - 1 = 2(u_n - 1) = 2v_n,
\]
sehingga $(v_n)$ \omterm{def:g11:seq:geometric}{geometri} dengan \omterm{def:g11:seq:geometric}{rasio} $2$ dan suku pertama
$v_0 = u_0 - 1 = 2$.

\emph{3.} Jadi $v_n = 2 \times 2^n = 2^{n+1}$ dan
$u_n = v_n + 1 = 2^{n+1} + 1$, yang membenarkan dugaan tadi. (Bilangan $1$
yang dikurangkan pada $v_n$ adalah \omterm{pb:g10:functions:1}{titik tetap} $x \mapsto 2x - 1$; gagasan
yang sama muncul kembali untuk $u_{n+1} = au_n + b$ pada \cref{ch:g12:seq}.)
\end{solution}

\begin{solution}{pb:g11:seq:1}
\textbf{1.} $h_1 = 1$, $h_2 = 3$, $h_3 = 7$.

\textbf{2.} Untuk memindahkan cakram terbesar, $n$ cakram di atasnya
harus lebih dahulu berpindah ke tiang cadangan ($h_n$ langkah); cakram besarnya
menyeberang ($1$ langkah); lalu $n$ cakram tadi harus mendaki kembali ke
atasnya ($h_n$ langkah): jadi $h_{n+1} = 2h_n + 1$.

\textbf{3.} $v_{n+1} = h_{n+1} + 1 = 2h_n + 2 = 2v_n$:
\omterm{def:g11:seq:geometric}{geometri} dengan \omterm{def:g11:seq:geometric}{rasio} $2$ dan $v_1 = 2$, sehingga $v_n = 2^n$ dan
$h_n = 2^n - 1$.

\textbf{4.} $2^{64} - 1 \approx 1.8 \times 10^{19}$ detik;
dibagi $3 \times 10^7$ detik per tahun, hasilnya
sekitar $6 \times 10^{11}$ tahun --- enam ratus miliar tahun,
empat puluh kali usia alam semesta. Para biarawan itu masih sempat
beristirahat minum kopi.

\textbf{5.} Pada setiap penyelesaian yang sah, tinjaulah langkah pertama
cakram terbawah: pada saat itu $n$ cakram lainnya harus semuanya berada di
satu tiang yang tersisa (paling sedikit $h_n$ langkah untuk membawanya ke
sana), dan setelah langkah terakhir cakram terbawah semuanya harus kembali
ke atasnya (paling sedikit $h_n$ langkah lagi): jadi setiap penyelesaian
menuntut paling sedikit $2h_n + 1$ langkah. Rekurensinya adalah lantai
sekaligus langit-langit: $2^n - 1$ memang optimum.

\textbf{6.} $1, 1, 2, 3, 5, 8, 13, 21, 34, 55, 89, 144$.

\textbf{7.} Bukan \omterm{def:g11:seq:arithmetic}{aritmetika} ($2 - 1 = 1$ tetapi $3 - 2 = 1$,
$5 - 3 = 2$: selisihnya berubah); bukan \omterm{def:g11:seq:geometric}{geometri}
($\frac21 = 2$ tetapi $\frac32 = 1.5$). Naik: untuk
$n \geq 2$ berlaku $F_{n+1} - F_n = F_{n-1} > 0$.

\textbf{8.} $F_k = F_{k+2} - F_{k+1}$, sehingga
\[
\sum_{k=1}^{n} F_k = (F_3 - F_2) + (F_4 - F_3) + \dots +
(F_{n+2} - F_{n+1}) = F_{n+2} - F_2 = F_{n+2} - 1 .
\]
Untuk $n = 6$: $1 + 1 + 2 + 3 + 5 + 8 = 20 = F_8 - 1 = 21 - 1$.

\textbf{9.} $F_k F_{k+1} - F_{k-1} F_k = F_k (F_{k+1} -
F_{k-1}) = F_k \cdot F_k = F_k^2$; menjumlahkannya secara teleskopik memberi
$F_n F_{n+1} - F_1 F_0$ (dengan $F_0 = 0$): jadi jumlah kuadratnya adalah
$F_n F_{n+1}$. Untuk $n = 4$: $1 + 1 + 4 + 9 = 15 = F_4 F_5 =
3 \times 5$.

\textbf{10.} $F_5 F_3 - F_4^2 = 5 \times 2 - 9 = 1$;
$F_6 F_4 - F_5^2 = 8 \times 3 - 25 = -1$;
$F_7 F_5 - F_6^2 = 13 \times 5 - 64 = 1$: berganti-ganti
$\pm 1$. Selisih satu antara $F_{n+1} F_{n-1}$ dan
$F_n^2$ inilah satu satuan persegi yang diperoleh atau hilang milik sang
pesulap: memotong persegi $F_n \times F_n$ menjadi potongan yang disusun
ulang menjadi persegi panjang $F_{n+1} \times F_{n-1}$ pasti menciptakan
atau menelan satu satuan --- yaitu celah tipis itu.

\textbf{11.} $F_{n+2} = F_{n+1} + F_n \geq F_n + F_n = 2F_n$
(\omterm{def:g11:seq:sequence}{barisannya} naik): setiap dua indeks, paling sedikit terjadi
pelipatduaan --- yaitu pertumbuhan paling sedikit setara \omterm{def:g11:seq:geometric}{geometri} berasio
$\sqrt2$ per indeks.

\textbf{12.} $1.5$; $1.667$; $1.6$; $1.625$; $1.615$;
$1.619$; $1.618$; $1.618$. Jika $r_n \to L$: dari
$F_{n+2} = F_{n+1} + F_n$, setelah dibagi $F_{n+1}$, diperoleh
$r_{n+1} = 1 + \frac{1}{r_n}$, sehingga $L = 1 + \frac1L$, yaitu
$L^2 = L + 1$: $L = \varphi = \frac{1 + \sqrt5}{2}$, yaitu \omterm{def:g11:seq:geometric}{rasio}
emas pada \cref{pb:g10:algebra:1}. Kelinci itu berkembang biak dalam
emas.

\textbf{13.} $v_{n+1} = u_{n+1} - \ell = a u_n + b - \ell$;
karena $\ell = a\ell + b$, bentuk itu sama dengan $a(u_n - \ell) = a v_n$:
jadi \omterm{def:g11:seq:geometric}{geometri} dengan \omterm{def:g11:seq:geometric}{rasio} $a$. Karena itu $v_n = a^n v_0$ dan
$u_n = a^n (u_0 - \ell) + \ell$.

\textbf{14.} \omterm{pb:g10:functions:1}{Titik tetapnya}: $\ell = 1.01\ell - 300$ memberi
$\ell = 30\,000$. Jadi
$d_n = 1.01^n (10\,000 - 30\,000) + 30\,000
= 30\,000 - 20\,000 \times 1.01^n$.

\textbf{15.} $d_n \leq 0$ menuntut $1.01^n \geq 1.5$:
$1.01^{40} \approx 1.489$, $1.01^{41} \approx 1.504$: jadi angsuran
ke-$41$ melunasi utangnya (dan sedikit lebih kecil daripada
$300$). Total angsurannya: sedikit di bawah $41 \times 300 = 12\,300$
euro --- jadi $10\,000$ yang dipinjam berongkos bunga sekitar
$2\,300$ euro.

\textbf{16.} \omterm{pb:g10:functions:1}{Titik tetapnya} $\ell = \frac{1000}{1 - 1.02} =
-50\,000$, sehingga $p_n = 1.02^n \times 100\,000 - 50\,000$. Setelah
$10$ tahun: $1.02^{10} \approx 1.219$, jadi
$p_{10} \approx 71\,900$ jiwa.

\textbf{17.} $\frac{1000 \times 1001}{2} = 500\,500$; lalu
$2^{20} - 1 = 1\,048\,575$.

\textbf{18.} Dari $7$ sampai $502$ dengan langkah $5$:
$\frac{502 - 7}{5} + 1 = 100$ suku; jumlahnya
$= 100 \times \frac{7 + 502}{2} = 25\,450$.

\textbf{19.} Saldonya $= 100 \times \frac{1.005^{60} - 1}
{1.005 - 1} \approx 100 \times \frac{0.3489}{0.005} \approx
6\,977$ euro --- yang $6\,000$ berasal dari setoran dan sekitar $977$
dari bunga: jumlah \omterm{def:g11:seq:geometric}{geometri} memang bahasa ibu bank.

\textbf{20.} Rumus eksplisit menjawab ``berapa $u_{1000}$''
seketika; rekurensi memerikan bagaimana sistemnya benar-benar berkembang ---
seninya adalah mengubah yang kedua menjadi yang pertama. \omterm{def:g11:seq:arithmetic}{Barisan aritmetika}
menambah, \omterm{def:g11:seq:geometric}{barisan geometri} mengalikan, dan setiap keluarga memiliki
rumus jumlahnya sendiri (pemasangan Gauss; kiat pelipatduaan). Kiat
\omterm{pb:g10:functions:1}{titik tetap} dan \omterm{def:g11:seq:sequence}{barisan} bantu mengubah setiap rekurensi afin
menjadi rekurensi \omterm{def:g11:seq:geometric}{geometri} --- pinjaman, populasi, dan
menara tadi semuanya jatuh olehnya. Fibonacci tidak menuruti kedua keluarga
itu, namun identitas teleskopik menangkap jumlah dan kuadratnya; potret
lengkapnya (rumus eksak, limit emas) masih menanti perkakas yang lebih kuat.

\end{solution}
